\documentclass[12pt,a4paper]{article}

% Compile with XeLaTeX for best font rendering.
\usepackage{fontspec}
\usepackage{polyglossia}
\setdefaultlanguage{russian}
\setotherlanguage{english}
\setmainfont{Inter}
\setsansfont{Inter}
\setmonofont{JetBrains Mono}

\usepackage[a4paper,margin=2.1cm]{geometry}
\usepackage{microtype}
\usepackage{setspace}
\setstretch{1.22}
\usepackage{parskip}
\setlength{\parindent}{0pt}

\usepackage{booktabs}
\usepackage{tabularx}
\usepackage{longtable}
\usepackage{array}
\usepackage{enumitem}
\setlist[itemize]{leftmargin=1.1em,itemsep=4pt,topsep=4pt}
\setlist[enumerate]{leftmargin=1.3em,itemsep=4pt,topsep=4pt}

\usepackage[table]{xcolor}
\definecolor{Navy}{HTML}{17324D}
\definecolor{Sky}{HTML}{2E86DE}
\definecolor{Emerald}{HTML}{1E9E66}
\definecolor{Amber}{HTML}{E0A106}
\definecolor{Crimson}{HTML}{C0392B}
\definecolor{Paper}{HTML}{F7FAFC}
\definecolor{Line}{HTML}{D9E2EC}
\definecolor{Muted}{HTML}{6B7785}

\usepackage{hyperref}
\usepackage{xurl}
\hypersetup{
  colorlinks=true,
  linkcolor=Sky,
  urlcolor=Sky,
  citecolor=Sky,
  pdftitle={MVP Audit Report -- New Level Hub},
  pdfauthor={New Level Hub Team}
}

\usepackage{titlesec}
\titleformat{\section}{\Large\bfseries\color{Navy}}{}{0em}{}
\titleformat{\subsection}{\large\bfseries\color{Sky}}{}{0em}{}
\titleformat{\subsubsection}{\normalsize\bfseries\color{Navy}}{}{0em}{}

\usepackage{fancyhdr}
\pagestyle{fancy}
\fancyhf{}
\lhead{\textcolor{Navy}{\textbf{New Level Hub}}}
\rhead{\textcolor{Muted}{MVP Audit}}
\cfoot{\textcolor{Muted}{\thepage}}
\renewcommand{\headrulewidth}{0.3pt}
\renewcommand{\headrule}{\hbox to\headwidth{\color{Line}\leaders\hrule height \headrulewidth\hfill}}

\usepackage[most]{tcolorbox}
\tcbset{
  arc=2mm,
  boxrule=0.4pt,
  colframe=Line,
  colback=Paper
}

\newtcolorbox{introbox}{
  colframe=Sky,
  colback=Sky!4
}
\newtcolorbox{riskbox}{
  colframe=Crimson,
  colback=Crimson!4
}

\newcommand{\done}{\textcolor{Emerald}{\textbf{[DONE]}}}
\newcommand{\partly}{\textcolor{Amber}{\textbf{[PARTIAL]}}}
\newcommand{\todoitem}{\textcolor{Crimson}{\textbf{[NOT DONE]}}}

\begin{document}
\sloppy

\begin{titlepage}
  \vspace*{1.6cm}
  {\Huge\bfseries\color{Navy} Итоговый статус реализации MVP\par}
  \vspace{0.2cm}
  {\LARGE\bfseries\color{Sky} New Level Hub\par}
  \vspace{1.2cm}

  \begin{tcolorbox}[colframe=Sky,colback=Sky!5]
    \textbf{Дата аудита:} 02.06.2026
  \end{tcolorbox}

  \vfill
  {\large\color{Muted} Technical documentation report}
\end{titlepage}

\tableofcontents
\newpage

\section{Как выполнен анализ}
\begin{tcolorbox}[colframe=Navy,colback=Paper,title=\textbf{Паспорт аудита},breakable]
\textbf{Период проведения:} 02.06.2026, примерно 13:13--13:31 (UTC+5).\\
\textbf{Исполнитель:} ИИ-агент Codex 5.3 (Cursor IDE).\\
\textbf{Формат проверки:} по чеклисту \texttt{mvp-audit.md}, с верификацией backend + frontend кодовой базы.\\
\textbf{Правило оценки:} \done = модель + API + UI; \partly = частично; \todoitem = не найдено.\\
\textbf{Исходный запрос:} ``Используй скилл \path|@NewLevelHub-Backend/.claude/agents/mvp-audit.md| и проведи полный аудит реализации платформы New Level Hub. Пройди по всем шагам скилла: 1) изучи структуру проекта, 2) пройди по каждому пункту чеклиста, 3) сформируй итоговый отчёт в указанном формате.''\\
\textbf{Референс-документ:} \path|@NewLevelHub-Backend/NewLevelHub-MVP.md|.
\end{tcolorbox}

\section{Итоговый статус реализации MVP — New Level Hub}
\textbf{Дата аудита:} 02.06.2026

\begin{introbox}
Аудит выполнен по вашему \texttt{mvp-audit.md}: сначала просмотрена структура backend/frontend, затем пройден чеклист блоков 1--7 с правилом оценки \done = модель + API + UI, \partly = частично, \todoitem = не найдено.
\end{introbox}

\subsection{Карта реализации (шаг 1)}
\begin{itemize}
  \item \textbf{Backend (Django apps):}
  \path|users|, \path|companies|, \path|bookings|, \path|crm|, \path|storage|, \path|hr|, \path|access|, \path|services|, \path|notifications|, \path|analytics|, \path|core|.\\
  Примеры опорных файлов: \path|NewLevelHub-Backend/apps/users/models.py|, \path|.../views.py|, \path|.../serializers.py|, \path|.../urls.py| в каждом модуле.
  \item \textbf{Frontend (pages/components/api):}
  Основной роутинг в \path|NewLevelHub-Web-Frontend/src/app/router.tsx|, API-слой в \path|src/shared/api/endpoints.ts| и \path|src/shared/api/client.ts|, страницы по блокам находятся в \path|src/pages/**|.
\end{itemize}

\subsection{Сводная таблица по блокам}
\rowcolors{2}{white}{Paper}
\begin{center}
\begin{tabularx}{\textwidth}{>{\raggedright\arraybackslash}X c c c c c}
\toprule
\textbf{Блок} & \textbf{Всего} & \textbf{\done} & \textbf{\partly} & \textbf{\todoitem} & \textbf{\% готовности} \\
\midrule
Блок 1. Пользователи и роли & 25 & 16 & 8 & 1 & 80\% \\
Блок 2. Компании & 23 & 17 & 6 & 0 & 87\% \\
Блок 3. Бронирование & 31 & 28 & 3 & 0 & 95\% \\
Блок 4. Корпоративный портал & 32 & 30 & 2 & 0 & 97\% \\
Блок 5. Доступ и пропуска & 14 & 12 & 2 & 0 & 93\% \\
Блок 6. Сервисы здания & 17 & 13 & 4 & 0 & 88\% \\
Блок 7. Уведомления и аналитика & 13 & 9 & 4 & 0 & 85\% \\
\midrule
\textbf{ИТОГО} & \textbf{155} & \textbf{125} & \textbf{29} & \textbf{1} & \textbf{90\%} \\
\bottomrule
\end{tabularx}
\end{center}

\textit{`\% готовности` рассчитан как \texttt{(DONE + 0.5 * PARTIAL) / N}.}

\section{Детализация по блокам}

\subsection{Блок 1. Система пользователей и ролей}
\textbf{\done}
\begin{itemize}
  \item Регистрация (включая invite-flow), логин, remember me, reset password (\path|apps/users/views.py|, \path|apps/users/serializers.py|, \path|src/pages/auth/**|).
  \item Email verification (\path|apps/users/models.py| \path|EmailVerificationToken|, \path|src/pages/auth/VerifyEmailPage.tsx|).
  \item Профиль: редактирование данных, аватар, роль/компания, настройки уведомлений (\path|src/pages/profile/ProfilePage.tsx|, \path|src/pages/notifications/NotificationPreferencesPage.tsx|).
  \item Роли и права (4 роли + role-based dashboard) (\path|apps/core/permissions.py|, \path|src/pages/dashboard/**|, \path|src/app/router.tsx|).
\end{itemize}
\textbf{\partly}
\begin{itemize}
  \item Создание админа компании суперадмином: через инвайты реализовано, отдельного явного UX ``создать админа компании'' нет (\path|apps/companies/...|, \path|src/pages/companies/CompanyMembersPage.tsx|).
  \item Ограничение доступа до подтверждения email: backend проверяет, frontend-guard неполный.
  \item Смена пароля: API есть, отдельной полноценной UI-формы в профиле нет.
  \item История действий пользователя: разрозненные данные есть, единого экрана/эндпоинта истории нет.
  \item Ролевые дашборды покрыты, но часть product-виджетов неполная.
\end{itemize}
\textbf{\todoitem}
\begin{itemize}
  \item Автовыход при неактивности (idle timeout) -- в коде не найден.
\end{itemize}

\subsection{Блок 2. Управление компаниями (мультитенанси)}
\textbf{\done}
\begin{itemize}
  \item CRUD компании, деактивация/активация, удаление с подтверждением (\path|apps/companies/views.py|, \path|src/pages/companies/CompanyListPage.tsx|).
  \item План/лимиты (employees/storage/boards) на уровне модели/API.
  \item Управление сотрудниками: список, инвайты, деактивация, удаление с переназначением задач (\path|apps/companies/views.py|, \path|src/pages/team/TeamManagePage.tsx|).
  \item Настройки компании: профиль, рабочие часы, инвайты, лимиты (\path|apps/companies/models.py| \path|CompanySettings|, \path|src/pages/companies/CompanySettingsPage.tsx|).
\end{itemize}
\textbf{\partly}
\begin{itemize}
  \item Лимиты тарифов в UI отображаются, но не в полном ``продуктовом'' виде по всем сценариям.
  \item Онбординг компании покрыт API/частично UI, но нет полного выделенного пошагового wizard-потока после создания компании.
  \item Usage-метрики лимитов на фронте частично разнесены по экранам.
\end{itemize}
\textbf{\todoitem}
\begin{itemize}
  \item Нет.
\end{itemize}

\subsection{Блок 3. Система бронирования ресурсов}
\textbf{\done}
\begin{itemize}
  \item Ресурсы, расписания, слоты, участники, описание, подтверждение, ``мои бронирования'' (\path|apps/bookings/**|, \path|src/pages/bookings/**|).
  \item Управление бронированиями (отмена, перенос, участники, admin/superadmin сценарии).
  \item Автоматизации: reminders, no-show, conflict-control, лимиты активных бронирований, recurring (\path|apps/bookings/tasks.py|, \path|apps/bookings/views.py|, \path|src/pages/bookings/RecurringBookingsPage.tsx|).
  \item Управление ресурсами суперадмином, включая bulk create (\path|src/pages/resources/**|).
\end{itemize}
\textbf{\partly}
\begin{itemize}
  \item Режим ``каталог: карточки/на карте'' реализован как близкие, но раздельные UX-потоки.
  \item Пункт ``появление в календаре'' работает через общую интеграцию календаря, но не как отдельный явный связанный UX-шаг.
  \item Отдельные продуктовые детали по капсулам/QR в UI неполные.
\end{itemize}
\textbf{\todoitem}
\begin{itemize}
  \item Нет.
\end{itemize}

\subsection{Блок 4. Корпоративный портал (CRM + workspace)}
\textbf{\done}
\begin{itemize}
  \item CRM Kanban почти полностью: доски, колонки, задачи, labels, drag\&drop, checklist, comments, attachments, history, filters, my tasks, archive (\path|apps/crm/**|, \path|src/pages/crm/**|).
  \item Справочник команды (\path|src/pages/team/TeamDirectoryPage.tsx|).
  \item Календарь команды (агрегация + day/week/month + фильтры) (\path|src/pages/calendar/CalendarPage.tsx|).
  \item Файловое хранилище: папки, upload/download, sharing, search/sort, usage/quota (\path|apps/storage/**|, \path|src/pages/storage/FileBrowserPage.tsx|).
  \item Leave requests + onboarding (\path|apps/hr/**|, \path|src/pages/hr/**|, \path|src/pages/onboarding/**|).
\end{itemize}
\textbf{\partly}
\begin{itemize}
  \item Клик по слоту календаря в действие ``создать бронь/задачу'' неполный.
  \item Attach ``из хранилища в задачу без дублирования'' поддержан в backend-модели, но frontend-flow не полностью явный.
\end{itemize}
\textbf{\todoitem}
\begin{itemize}
  \item Нет.
\end{itemize}

\subsection{Блок 5. Система доступа и гостевые пропуска}
\textbf{\done}
\begin{itemize}
  \item CRUD пропусков, QR-валидация, revoke/resend, журналы доступа, фильтры, CSV export (\path|apps/access/**|, \path|src/pages/access/**|).
  \item Веб-валидация для ресепшн (\path|src/pages/access/PassValidatePage.tsx|).
\end{itemize}
\textbf{\partly}
\begin{itemize}
  \item ``Любой сотрудник/гость коворкинга'' -- для guest роль ограничена.
  \item Отправка QR при создании покрыта не полностью как guaranteed UX-цепочка (основной акцент на resend/detail).
\end{itemize}
\textbf{\todoitem}
\begin{itemize}
  \item Нет.
\end{itemize}

\subsection{Блок 6. Сервисы здания}
\textbf{\done}
\begin{itemize}
  \item Карта этажей, точки, поиск, floor switching, CRUD точек/этажей (\path|apps/services/**|, \path|src/pages/map/**|).
  \item Сервисные заявки: типы, статусы, assign, rating, quick cleaning (\path|src/pages/service-requests/**|).
  \item Объявления БЦ/компаний, read-state (\path|apps/services/models.py| \path|AnnouncementRead|, \path|src/pages/announcements/**|).
\end{itemize}
\textbf{\partly}
\begin{itemize}
  \item Клик по ресурсу карты в прямой booking-flow неполный.
  \item Статусная визуализация точек (по всем продуктовым ожиданиям) частично.
  \item \path|MapManagePage| на фронте заглушка; часть управления перенесена в основной map UI.
  \item Email-рассылка важных объявлений есть в backend, frontend-контролы частичные.
\end{itemize}
\textbf{\todoitem}
\begin{itemize}
  \item Нет.
\end{itemize}

\subsection{Блок 7. Уведомления и аналитика}
\textbf{\done}
\begin{itemize}
  \item Центр уведомлений, read/read-all, preferences, DND (\path|apps/notifications/**|, \path|src/pages/notifications/**|).
  \item Аналитика superadmin/company admin + CSV export (\path|apps/analytics/views.py|, \path|src/pages/analytics/**|).
\end{itemize}
\textbf{\partly}
\begin{itemize}
  \item ``Колокольчик + бейдж'' и часть UX-интеграции в глобальном хедере покрыты неполно.
  \item Heatmap пиковых часов выражена неполно в UI.
  \item Export в PDF не найден (есть CSV).
  \item HTML email templates как полноценный шаблонный пакет явно не прослеживаются end-to-end.
\end{itemize}
\textbf{\todoitem}
\begin{itemize}
  \item Нет.
\end{itemize}

\section{Ключевые риски и наблюдения}
\begin{riskbox}
\begin{itemize}
  \item \textbf{Единственный явный [NOT DONE]:} idle auto-logout отсутствует -- это gap безопасности и соответствия продукту.
  \item \textbf{Критичный UX-долг:} ряд backend-фич уже есть, но frontend-слой частично не доведен (password change UI, action history, map manage UX, calendar slot actions).
  \item \textbf{Роутинг/доступность экранов:} есть реализованные экраны, не полностью подключенные в навигации (например лог доступа).
  \item \textbf{Аналитика:} функционально сильная, но не закрывает полный продуктовый набор визуализаций и PDF export.
\end{itemize}
\end{riskbox}

\section{Рекомендуемые следующие шаги}
\begin{enumerate}
  \item Закрыть Block 1 технический долг: \texttt{idle timeout}, полноценный \texttt{change password} UI, единый \texttt{user activity history}.
  \item Довести frontend-паритет с backend в блоках 5--7: map manage, access log routing, notification bell UX.
  \item Закрыть аналитические gaps: heatmap и PDF export.
  \item Для ``PARTIAL'' пунктов ввести Definition of Done ``модель + API + UI + e2e smoke test'' и пройтись итеративно блоками 1$\rightarrow$7.
  \item После доработок сделать повторный аудит по тому же шаблону и зафиксировать delta-отчет.
\end{enumerate}

\end{document}
